\section{[30pt] System}
Consider the unicycle 4-state model. The state $x = [p_x,p_y,v,\theta]^T\in\mathbb{R}^4$ contains the 2D position $(p_x,p_y)$, the speed $v$, and the heading $\theta$ of the vehicle. The control input $u=[\dot v, \dot \theta]^T\in\mathbb{R}^2$ contains the acceleration $\dot v$ and the turning rate $\dot\theta$. The dynamic equation is
\begin{eqnarray}\label{eq: dynamics}
\dot x = \left[\begin{array}{c}
v\cos\theta\\
v\sin\theta\\
0\\
0
\end{array} \right] + \left[\begin{array}{cc}
0 & 0\\
0 & 0\\
1 & 0\\
0 & 1
\end{array} \right] u
\end{eqnarray}

{1.1 [1pt] Is this system control-affine?}

\solution{Yes, the system is control affine because the control input shows up linearly in the dynamics.}

{1.2 [5pt] Which description is most appropriate? Choose one from each row}
\begin{center}
    \begin{tabular}{cc}
Linear &   Nonlinear\\
Continuous-Time & Discrete-Time\\
Time-invariant & Time-varying\\
Deterministic & Stochastic\\
Continuous-State & Discrete-State
\end{tabular}
\end{center}

\solution{The following description is appropriate:
\begin{itemize}
    \item Nonlinear
    \item Continuous-Time
    \item Time-invariant
    \item Deterministic
    \item Continuous-State
\end{itemize}}

% \grading{The following description is appropriate:
\begin{itemize}
    \item Nonlinear
    \item Continuous-Time
    \item Time-invariant
    \item Deterministic
    \item Continuous-State
\end{itemize}}

{1.3 [8pt] Linearize the system at $[p_x^r,p_y^r,v^r,\theta^r]^T$. Write down the resulting dynamics.}

\solution{The resulting dynamics are defined by
\begin{equation}
    \delta \dot{x} = 
    \begin{bmatrix}
        0 & 0 & \cos\theta^r & -v\sin\theta^r \\
        0 & 0 & \sin\theta^r & v\cos\theta^r \\
        0 & 0 & 0 & 0 \\
        0 & 0 & 0 & 0
    \end{bmatrix} \delta x + 
    \begin{bmatrix}
        0 & 0 \\
        0 & 0 \\
        1 & 0 \\
        0 & 1
    \end{bmatrix} \delta u
\end{equation}
where $\delta x = x - x^r$, $x^r = [p_x^r,p_y^r,v^r,\theta^r]^T$, $\delta u = u - u^r$.}

% \grading{The resulting dynamics are defined by
\begin{equation}
    \delta \dot{x} = 
    \begin{bmatrix}
        0 & 0 & \cos\theta^r & -v\sin\theta^r \\
        0 & 0 & \sin\theta^r & v\cos\theta^r \\
        0 & 0 & 0 & 0 \\
        0 & 0 & 0 & 0
    \end{bmatrix} \delta x + 
    \begin{bmatrix}
        0 & 0 \\
        0 & 0 \\
        1 & 0 \\
        0 & 1
    \end{bmatrix} \delta u
\end{equation}
where $\delta x = x - x^r$, $x^r = [p_x^r,p_y^r,v^r,\theta^r]^T$, $\delta u = u - u^r$.}

{1.4 [8pt] Compute the linearized system dynamics at $[1,1,0,0]^T$. Is the linearized system at this point controllable? Justify your answer.}

\solution{The linearized dynamics at $[1,1,0,0]^T$ are
\begin{align}
    \delta \dot{x} &= 
    \begin{bmatrix}
        0 & 0 & \cos\theta^r & -v\sin\theta^r \\
        0 & 0 & \sin\theta^r & v\cos\theta^r \\
        0 & 0 & 0 & 0 \\
        0 & 0 & 0 & 0
    \end{bmatrix} \delta x + 
    \begin{bmatrix}
        0 & 0 \\
        0 & 0 \\
        1 & 0 \\
        0 & 1
    \end{bmatrix} \delta u \\
    &= \begin{bmatrix}
        0 & 0 & 1 & 0 \\
        0 & 0 & 0 & 0 \\
        0 & 0 & 0 & 0 \\
        0 & 0 & 0 & 0
    \end{bmatrix} \delta x + 
    \begin{bmatrix}
        0 & 0 \\
        0 & 0 \\
        1 & 0 \\
        0 & 1
    \end{bmatrix} \delta u.
\end{align}
This linearized system is not controllable because we have no lever to affect the $p_y$ state.}

% \grading{The linearized dynamics at $[1,1,0,0]^T$ are
\begin{align}
    \delta \dot{x} &= 
    \begin{bmatrix}
        0 & 0 & \cos\theta^r & -v\sin\theta^r \\
        0 & 0 & \sin\theta^r & v\cos\theta^r \\
        0 & 0 & 0 & 0 \\
        0 & 0 & 0 & 0
    \end{bmatrix} \delta x + 
    \begin{bmatrix}
        0 & 0 \\
        0 & 0 \\
        1 & 0 \\
        0 & 1
    \end{bmatrix} \delta u \\
    &= \begin{bmatrix}
        0 & 0 & 1 & 0 \\
        0 & 0 & 0 & 0 \\
        0 & 0 & 0 & 0 \\
        0 & 0 & 0 & 0
    \end{bmatrix} \delta x + 
    \begin{bmatrix}
        0 & 0 \\
        0 & 0 \\
        1 & 0 \\
        0 & 1
    \end{bmatrix} \delta u.
\end{align}
This linearized system is not controllable because we have no lever to affect the $p_y$ state.}

{1.5 [8pt] Under which condition on the parameters $p_x^r, p_y^r, v^r, \theta^r$ is the linearized system controllable? Find sufficient and necessary conditions.}

\solution{The linearized system is controllable if and only if
$\operatorname{rank}(\mathcal{C}) = n$, where $\mathcal{C}$ is the
controllability matrix and $n$ is the number of the dimensions. Since $n=4$ for
the unicycle model, we know the controllability matrix is
\begin{equation*}
    \mathcal{C} = \left[ A^3B ~ A^2B ~ AB ~ B \right].
\end{equation*} 
Evaluating this expression, we arrive at
\begin{equation*}
    \mathcal{C} = 
    \left[\begin{array}{cc|cc|cc|cc}
        0 & 0 & 0 & 0 & \cos\theta^r & -v^r\sin\theta^r & 0 & 0 \\
        0 & 0 & 0 & 0 & \sin\theta^r & v^r\cos\theta^r  & 0 & 0 \\
        0 & 0 & 0 & 0 & 0 & 0 & 1 & 0 \\
        0 & 0 & 0 & 0 & 0 & 0 & 0 & 1
    \end{array}\right].
\end{equation*}
Since the first 4 columns are zero, we can investigate the "right" $4 \times 4$
square matrix. From the 2 bottom rows, we see they are immediately linearly
independent, giving an immediate $\operatorname{rank}(\mathcal{C}) \geq 2$.
Thus, we can further simplify the problem and investigate the rank of 
\begin{equation*}
    \mathcal{M} = 
    \left[\begin{array}{cc}
        \cos\theta^r & -v^r\sin\theta^r \\
        \sin\theta^r & v^r\cos\theta^r 
    \end{array}\right]
\end{equation*}
If the determinant of $\mathcal{M}$ is non-zero, then $\mathcal{M}$ itself is
full rank, and $\operatorname{rank}(\mathcal{C}) = 4$. Thus, we evaluate 
\begin{align*}
    \det({\mathcal{M}}) &= v^r\cos^2\theta^r + v^r\sin^2\theta^r \\
    &= v^r(\cos^2\theta^r + \sin^2\theta^r) \\
    &= v^r
\end{align*}
Thus, $\det({\mathcal{M}}) \neq 0$ if and only if $v^r \neq 0$, and therefore,
the system is controllable if and only if $v^r \neq 0$. }

% \grading{The linearized system is controllable if and only if
$\operatorname{rank}(\mathcal{C}) = n$, where $\mathcal{C}$ is the
controllability matrix and $n$ is the number of the dimensions. Since $n=4$ for
the unicycle model, we know the controllability matrix is
\begin{equation*}
    \mathcal{C} = \left[ A^3B ~ A^2B ~ AB ~ B \right].
\end{equation*} 
Evaluating this expression, we arrive at
\begin{equation*}
    \mathcal{C} = 
    \left[\begin{array}{cc|cc|cc|cc}
        0 & 0 & 0 & 0 & \cos\theta^r & -v^r\sin\theta^r & 0 & 0 \\
        0 & 0 & 0 & 0 & \sin\theta^r & v^r\cos\theta^r  & 0 & 0 \\
        0 & 0 & 0 & 0 & 0 & 0 & 1 & 0 \\
        0 & 0 & 0 & 0 & 0 & 0 & 0 & 1
    \end{array}\right].
\end{equation*}
Since the first 4 columns are zero, we can investigate the "right" $4 \times 4$
square matrix. From the 2 bottom rows, we see they are immediately linearly
independent, giving an immediate $\operatorname{rank}(\mathcal{C}) \geq 2$.
Thus, we can further simplify the problem and investigate the rank of 
\begin{equation*}
    \mathcal{M} = 
    \left[\begin{array}{cc}
        \cos\theta^r & -v^r\sin\theta^r \\
        \sin\theta^r & v^r\cos\theta^r 
    \end{array}\right]
\end{equation*}
If the determinant of $\mathcal{M}$ is non-zero, then $\mathcal{M}$ itself is
full rank, and $\operatorname{rank}(\mathcal{C}) = 4$. Thus, we evaluate 
\begin{align*}
    \det({\mathcal{M}}) &= v^r\cos^2\theta^r + v^r\sin^2\theta^r \\
    &= v^r(\cos^2\theta^r + \sin^2\theta^r) \\
    &= v^r
\end{align*}
Thus, $\det({\mathcal{M}}) \neq 0$ if and only if $v^r \neq 0$, and therefore,
the system is controllable if and only if $v^r \neq 0$. }


\section{[46pt] Model Mismatch}

In this section, the controller imagines that the vehicle follows the unicycle
4-state model

\begin{equation}
x_u=
\begin{bmatrix}p_x\\p_y\\v\\\theta\end{bmatrix},
\qquad
u_u=
\begin{bmatrix}\dot v\\\dot\theta\end{bmatrix},
\end{equation}

where $(p_x,p_y)$ is the position in the fixed world frame, $v$ is the forward
speed, $\theta$ is the heading, $\dot v$ is the longitudinal acceleration, and
$\dot\theta$ is the turning rate. The actual rollout uses the dynamic bicycle
model supplied in \texttt{hw1/vehicle\_models.py}. Both transitions use forward
Euler integration with time step $\Delta t=0.01\,\mathrm{s}$. Unless otherwise
stated, use

\begin{equation}
x_u(0)=
\begin{bmatrix}0\\0\\0.5\\0\end{bmatrix},
\qquad
p_g=
\begin{bmatrix}5\\5\end{bmatrix},
\qquad
T=10\,\mathrm{s}.
\end{equation}

Here, $x_u(0)$ is the initial unicycle state, $p_g$ is the goal position with
components $p_{g,x}$ and $p_{g,y}$, and $T$ is the rollout duration.

\paragraph{2.1 [15pt] Unicycle feedback controller.}
Define the goal-position error $e_p\in\mathbb{R}^2$, distance to the goal
$\rho\in\mathbb{R}_{\geq 0}$, desired heading $\theta_d$, and wrapped heading
error $e_\theta$ by

\begin{align}
e_p & =
\begin{bmatrix}p_{g,x}-p_x\\p_{g,y}-p_y\end{bmatrix},
&
\rho & =\lVert e_p\rVert_2,\\
\theta_d & =\tan^{-1}\!\left(\frac{e_{p,y}}{e_{p,x}}\right),
&
e_\theta & =\operatorname{wrap}(\theta_d-\theta)\in[-\pi,\pi).
\end{align}

Here, $\theta_d$ is the angle of the position-error vector measured from the
positive $x$-axis. The operator $\operatorname{wrap}(\phi)$ maps an angle
$\phi$ to $[-\pi,\pi)$. In code, compute this operation with the provided
\texttt{wrap\_angle()} function; do not implement a separate angle-wrapping
rule.

Let $k_v>0$ be the velocity-feedback gain, $k_\theta>0$ be the
heading-feedback gain, $v_{\max}>0$ be the maximum desired speed, and $v_d$ be
the desired speed defined by

\begin{equation}
v_d=\min(v_{\max},\rho)\max(0,\cos e_\theta),
\end{equation}

Implement the two controller calculations

\begin{equation}\label{eq:unicycle-controller}
u_u=
\begin{bmatrix}\dot v\\\dot\theta\end{bmatrix}
=
\begin{bmatrix}
k_v(v_d-v)\\
k_\theta e_\theta
\end{bmatrix}.
\end{equation}

Complete \texttt{unicycle\_control} in \texttt{problem.py}. Use the default
gains $(k_v,k_\theta)=(2.0,2.5)$, run the unicycle model, and submit the XY
trajectory and every state-versus-time plot. Explain the purpose of each gain
and of the factor $\max(0,\cos e_\theta)$. The supplied return lines already
clip the acceleration and turning rate. After computing the errors and desired
speed, you only need to calculate \texttt{v\_dot} and
\texttt{theta\_dot}; do not add another saturation operation.

% \solution{\input{soln/2.1_26fall}}
% \grading{\input{grad/2.1_26fall}}

\paragraph{2.2 [15pt] Pure-rolling bicycle command conversion.}
The dynamic bicycle state and control follow the Lecture 2 and SPARK order

\begin{equation}
x_b=
\begin{bmatrix}X\\Y\\v_x\\v_y\\r\\\psi\end{bmatrix},
\qquad
u_b=
\begin{bmatrix}a_x\\\delta\end{bmatrix},
\end{equation}

where $(X,Y)$ is the position in the fixed world frame, $v_x$ and $v_y$ are the
longitudinal and lateral body-frame velocities, $r=\dot\psi$ is the yaw rate,
$\psi$ is the heading, $a_x$ is the longitudinal acceleration command, and
$\delta$ is the front-wheel steering angle. Implement the initial-state
conversion

\begin{equation}\label{eq:state-conversion}
x_b(0)=
\begin{bmatrix}
p_x(0)\\p_y(0)\\v(0)\\0\\0\\\theta(0)
\end{bmatrix}.
\end{equation}

At each step, evaluate the Problem 2.1 controller using the virtual unicycle
state $\widetilde{x}_u$, defined by

\begin{equation}
\widetilde{x}_u=
\begin{bmatrix}X\\Y\\v_x\\\psi\end{bmatrix}.
\end{equation}

The actual rollout plant remains the six-state SPARK dynamic bicycle, whose
lateral motion is generated by front and rear tire forces. For the purpose of
converting the unicycle command only, replace those tire-force equations by
the no-side-slip, pure-rolling approximation. Here, $\alpha_f$ and $\alpha_r$
denote the front and rear tire slip angles, and $l_f>0$ and $l_r>0$ denote the
distances from the vehicle center of mass to the front and rear axles,
respectively. The pure-rolling constraints are

\begin{equation}
\alpha_f=\alpha_r=0
\quad\Longrightarrow\quad
v_y-l_r r=0,
\qquad
v_y+l_f r=v_x\tan\delta.
\end{equation}

With wheelbase $L=l_f+l_r$, these constraints give

\begin{equation}
r=\dot\psi=\frac{v_x}{L}\tan\delta.
\end{equation}

Therefore, if the shared controller returns $(\dot v,\dot\theta)$, implement
the provided command conversion

\begin{equation}\label{eq:control-conversion}
u_b=
\begin{bmatrix}a_x\\\delta\end{bmatrix}
=
\begin{bmatrix}
\dot v\\
\operatorname{sat}_{\delta_{\max}}
\!\left(\tan^{-1}\!\left(\frac{L\dot\theta}{v_x}\right)\right)
\end{bmatrix}.
\end{equation}

Here, $\delta_{\max}>0$ is the steering-angle limit and
$\operatorname{sat}_{\delta_{\max}}(z)=
\min(\delta_{\max},\max(-\delta_{\max},z))$.

The inverse tangent is taken on its principal branch, and the conversion is
used on the forward-motion domain $v_x>0$. The acceleration command already
satisfies its limit from Problem 2.1. Use the pure-rolling relation only for
the command conversion: do not replace the supplied tire-force bicycle
transition used for the actual rollout.

Complete \texttt{bicycle\_initial\_state} and \texttt{bicycle\_control}, run
the bicycle model, and submit the XY trajectory and every state-versus-time
plot. Explain the tracking error caused by model mismatch. Hint: relate the
error to the bicycle's lateral velocity $v_y$ and yaw rate $r$.

% \solution{\input{soln/2.2_26fall}}
% \grading{\input{grad/2.2_26fall}}

\paragraph{2.3 [16pt] Two-parameter performance analysis.}
Use the shared controller for both models and evaluate all nine gain pairs in

\begin{equation}
k_v\in\{1.5,2.0,2.5\},
\qquad
k_\theta\in\{1.5,2.5,3.5\}.
\end{equation}

At matching sample indices $k$, corresponding to time $t_k=k\Delta t$, define
the unicycle--bicycle position mismatch $e_p^{\mathrm{mm}}(k)$ and orientation
mismatch $e_\theta^{\mathrm{mm}}(k)$ by

\begin{equation}
e_p^{\mathrm{mm}}(k)=\left\|
\begin{bmatrix}p_{x,k}\\p_{y,k}\end{bmatrix}
-
\begin{bmatrix}X_k\\Y_k\end{bmatrix}
\right\|_2,
\qquad
e_\theta^{\mathrm{mm}}(k)=
\left|\operatorname{wrap}(\theta_k-\psi_k)\right|.
\end{equation}

Here, $(p_{x,k},p_{y,k},\theta_k)$ are components of the unicycle state at
$t_k$, while $(X_k,Y_k,\psi_k)$ are components of the bicycle state at the
same time. Use the provided \texttt{wrap\_angle()} function for the angular
difference.

Run the provided gain study and submit exactly two figures: (i) position error
versus time and (ii) orientation error versus time. Each figure should show all
nine gain pairs with a readable legend. Analyze how $k_v$ and $k_\theta$ affect
the transient and final errors, and discuss whether one gain pair is uniformly
best for both errors.

% \solution{\input{soln/2.3_26fall}}
% \grading{\input{grad/2.3_26fall}}

\section{[24pt] True or False}

3.1 [3pt] There is a unique equilibrium point for any system.

% \solution{\input{soln/3.1}}

3.2 [3pt] For linear systems, local stability implies global stability.

% \solution{\input{soln/3.2}}

3.3 [3pt] A stable system is bounded.

% \solution{\input{soln/3.3}}

3.4 [3pt] A stable system converges to its equilibrium.

% \solution{\input{soln/3.4}}

3.5 [3pt] Adding additional actuators to a system can always improve its controllability.

% \solution{\input{soln/3.5}}

3.6 [3pt] Unobservability occurs due to system stochasticity.

% \solution{\input{soln/3.6}}

3.7 [3pt] For a Markov process, the future state depends on the past state even when the current state and control are given.

% \solution{\input{soln/3.7}}

3.8 [3pt] For any nonlinear system with known model and known initial state, the future can be fully predicted by simulation.

% \solution{\input{soln/3.8}}
